% Cascade, Cortex, JetSki: A Three-Layer Telemetry Anatomy
% of a Commercial AI Coding Agent
%
% Target: ACM AISec 2026 (CCS Workshop, The Hague, 11/15)
% Format: ACM double-column, ≤6 pages + ≤2 pages refs/appendix
% Status: §1–§6 complete
%
% Anonymization: Cascade/Cortex/JetSki are internal codenames
% extracted from telemetry identifiers. No vendor names appear.

\documentclass[sigconf,anonymous]{acmart}

%%% ACM metadata (placeholder for submission)
\acmConference[AISec '26]{Workshop on Artificial Intelligence and Security}{November 15, 2026}{The Hague, Netherlands}
\acmYear{2026}
\acmISBN{}
\acmDOI{}

%%% Anonymization macros — single point of control
\newcommand{\cascade}{\textsc{Cascade}\xspace}
\newcommand{\cortex}{\textsc{Cortex}\xspace}
\newcommand{\jetski}{\textsc{JetSki}\xspace}
\newcommand{\battlestar}{\textsc{Battlestar}\xspace}
\newcommand{\agentname}{the analyzed agent\xspace}  % generic reference
\newcommand{\uss}{\textsc{USS}\xspace}              % Unified State Sync

%%% Packages
\usepackage{xspace}
\usepackage{booktabs}
\usepackage{enumitem}
\usepackage{balance}

\begin{document}

\title{Cascade, Cortex, JetSki: A Three-Layer Telemetry Anatomy of a Commercial AI Coding Agent}

\begin{abstract}
AI coding agents---IDE-integrated systems that combine large language
models with tool-execution capabilities---are rapidly becoming standard
development infrastructure, operating with deep access to file systems,
terminals, browsers, and version-control history.  Yet no systematic
study has examined the telemetry architectures through which these
agents report developer activity to their vendors.

We present the first structural anatomy of such a telemetry
architecture, derived from protobuf schemas extracted from a shipping
binary.  The architecture comprises three subsystems---\cascade (UI),
\cortex (engine), and \jetski (backend)---defining \textbf{90+ step
types}, \textbf{53 reverse-gRPC RPCs}, \textbf{14 IDE identifiers},
and \textbf{7 metric categories} including character-level edit
attribution.  We trace cross-layer data flow from keystroke to cloud,
argue per RFC~7258 that the pipeline constitutes \emph{pervasive
monitoring}, present a cross-agent comparison revealing extreme
\emph{control asymmetry} (1 user-accessible control vs.\ 324
telemetry fields), and propose a five-stage behavioral deviation
model grounded in observed agent misalignment episodes.  We close by
contrasting this architecture with data-minimization principles from
self-auditing security systems.
\end{abstract}

\begin{CCSXML}
<ccs2012>
   <concept>
       <concept_id>10002978.10002991.10002994</concept_id>
       <concept_desc>Security and privacy~Privacy protections</concept_desc>
       <concept_significance>500</concept_significance>
   </concept>
   <concept>
       <concept_id>10002978.10002979.10002983</concept_id>
       <concept_desc>Security and privacy~Intrusion detection systems</concept_desc>
       <concept_significance>300</concept_significance>
   </concept>
   <concept>
       <concept_id>10002978.10003014.10003017</concept_id>
       <concept_desc>Security and privacy~Software reverse engineering</concept_desc>
       <concept_significance>500</concept_significance>
   </concept>
   <concept>
       <concept_id>10010147.10010178.10010179</concept_id>
       <concept_desc>Computing methodologies~Autonomous agents</concept_desc>
       <concept_significance>500</concept_significance>
   </concept>
   <concept>
       <concept_id>10002978.10002991.10002995</concept_id>
       <concept_desc>Security and privacy~Pseudonymity, anonymity and untraceability</concept_desc>
       <concept_significance>300</concept_significance>
   </concept>
</ccs2012>
\end{CCSXML}

\ccsdesc[500]{Security and privacy~Privacy protections}
\ccsdesc[300]{Security and privacy~Intrusion detection systems}
\ccsdesc[500]{Security and privacy~Software reverse engineering}
\ccsdesc[500]{Computing methodologies~Autonomous agents}
\ccsdesc[300]{Security and privacy~Pseudonymity, anonymity and untraceability}

\keywords{AI agent telemetry, pervasive monitoring, protobuf schema extraction, developer privacy, control asymmetry, RFC~7258, reverse-gRPC}

\maketitle

%% ================================================================
%%  §1  INTRODUCTION
%% ================================================================
\section{Introduction}\label{sec:intro}

The adoption of AI coding agents has accelerated dramatically since
2024.  Commercial offerings now ship as IDE extensions or standalone
tools that pair a large language model (LLM) with an \emph{agentic
runtime}: an orchestration layer capable of reading files, executing
terminal commands, invoking browser automation, and spawning
sub-agents---all within the developer's local environment.

This deep integration raises a question that has received remarkably
little scrutiny: \emph{what telemetry do these agents collect, and
where does it go?}  RFC~7258~\cite{rfc7258} (BCP~188) established
that ``pervasive monitoring is an attack,'' defined by \emph{scale},
not purpose.  An AI agent that records code content, developer intent,
execution context, and workspace structure as a structured trajectory
streamed to a cloud backend produces a full-stack developer activity
record with no precedent in software development tooling.

In this paper, we present the first public, structural anatomy of such
a telemetry architecture.  Our analysis targets a single, commercially
deployed AI coding agent whose internal architecture is organized
around three named subsystems:

\begin{description}[style=unboxed,leftmargin=0pt,nosep]
  \item[\cascade] The UI front-end layer, implemented as an IDE
    extension.  It captures 78~distinct UI event types and exposes
    53~reverse-gRPC RPCs through which the agent engine can
    \emph{push} operations to the IDE.
  \item[\cortex] The AI agent engine, implemented as a Go-based
    Language Server.  It defines 90+~step types (from file reads to
    browser pixel-clicks), a 12-state execution state machine, and
    14~trajectory types---collectively forming a complete behavioral
    log of every action the agent takes.
  \item[\jetski] The backend bridge, connecting the local agent to
    a cloud-hosted internal API service.\footnote{We retain the internal
    API namespace \texttt{\{vendor\}.internal.\{service\}.v1internal} as
    extracted from protobuf definitions; the actual vendor and service
    identifiers are redacted to comply with double-blind review.  This
    is a structural observation, not a vendor identification.}  It streams incremental
    state updates, aggregates metrics to analytics infrastructure, and
    manages user onboarding and tier classification across 14~registered
    IDE types.
\end{description}

\noindent Our analysis is based entirely on protobuf schema definitions
extracted from shipping binaries (22~\texttt{.proto} files,
$\sim$360\,KB of schema definitions) and validated against
100~decrypted protobuf binary dumps of actual trajectory data.  We do
not reverse-engineer server-side systems, intercept network traffic, or
decompile application logic.  All codenames (\cascade, \cortex,
\jetski) are the system's own internal identifiers, extracted from
protobuf package declarations and telemetry event prefixes.\footnote{All
proto definitions and supporting evidence are available at
\url{https://anonymous.4open.science/r/sass-evidence-prerevirew-177fj2A}.}

\paragraph{Contributions.}
\begin{enumerate}[nosep]
  \item The first public structural anatomy of a three-layer telemetry
    architecture in a commercial AI coding agent (\S\ref{sec:anatomy}).
  \item Quantification of the telemetry surface area: 90+~step types,
    53~reverse-gRPC RPCs, 14~IDE type identifiers, 7+1~metric event
    categories, 22~proto files (\S\ref{sec:anatomy}).
  \item A cross-agent comparison of four commercial agents, revealing
    extreme control asymmetry (\S\ref{sec:comparison}).
  \item A five-stage behavioral deviation progression model (S1--S5)
    grounded in observed agent misalignment, with corresponding
    defense responses (\S\ref{sec:deviation-model}).
  \item Discussion of implications for defense system design, with
    reference to data-minimization principles as a constructive
    counterexample (\S\ref{sec:discussion}).
\end{enumerate}

\paragraph{Relationship to companion work.}
A companion paper~\cite{rogue-agent}, submitted to the same venue,
documents three independent incidents of autonomous agent behavioral
divergence in a production environment, presenting 3.2\,MB of forensic
evidence including the agent's own written outputs, and proposes
SASS---a transport-layer containment protocol.  An ongoing companion
investigation~\cite{ring0ring3} presents preliminary evidence suggesting
that the \texttt{PlanningModeConfig} propagated through
the USS dual-use channel (\S\ref{sec:cascade}) determines system
prompt assembly---meaning the telemetry infrastructure documented here
may double as the delivery mechanism for prompt priority overrides.  The
present paper focuses exclusively on the \emph{telemetry architecture};
the three papers share no overlapping sections and cite each other only
for context.


%% ================================================================
%%  §2  BACKGROUND
%% ================================================================
\section{Background}\label{sec:background}

\subsection{RFC~7258 and Telemetry Evolution}
\label{sec:rfc7258}

RFC~7258~\cite{rfc7258} (BCP~188, 2014) codified the position that
``pervasive monitoring is a technical attack,'' motivated by the
Snowden disclosures~\cite{snowden2013}.  Three aspects are relevant:
(1)~\emph{scale} as the defining criterion---breadth, not identity of
the monitor, determines pervasiveness;
(2)~\emph{infrastructure as attack surface}---collected data ``can be
subpoenaed, stolen, or repurposed'';
(3)~\emph{design-level mitigation}---controls must be architectural,
not merely operational.  RFC~7258 influenced TLS~1.3~\cite{tls13},
DNS-over-HTTPS~\cite{doh}, and QUIC~\cite{quic}; we extend it to
AI agent telemetry.

Development-tool telemetry has evolved through three phases:
crash reporting (pre-2015; stack traces, OS version),
feature-usage analytics (2015--2023; commands invoked, panels
opened~\cite{vscode-telemetry,jetbrains-fus}), and
AI agent telemetry (2024--present).  Phase~3 is qualitatively
different: agents observe \emph{code content}, \emph{developer intent}
(natural-language prompts), \emph{execution context} (terminal output,
browser DOM), and \emph{workspace structure} (git history, directory
listings).  When recorded as a structured trajectory and streamed to
a cloud backend, the result is a machine-readable diary of the
developer's work session.

\subsection{Architecture Patterns and Threat Model}
\label{sec:agent-arch}

Modern AI coding agents converge on a three-tier architecture:
(1)~a \textbf{UI Layer} (IDE extension capturing input and rendering
output), (2)~an \textbf{Agent Engine} (LLM orchestration runtime
managing prompt assembly, tool dispatch, and trajectory management),
and (3)~a \textbf{Backend Service} (cloud API for inference,
persistence, and telemetry aggregation).  Each layer independently
generates, enriches, and transmits telemetry.  Prior work has examined
telemetry in browsers~\cite{browser-telemetry}, mobile
apps~\cite{mobile-telemetry}, and OSes~\cite{os-telemetry}, but no
study has mapped an AI coding agent's telemetry architecture.

We consider a passive adversary with access to the telemetry stream
at any pipeline point: the vendor itself, a state-level actor with
compulsory disclosure authority, or an infrastructure attacker who has
compromised the backend.  Following RFC~7258, we treat the pipeline's
\emph{existence} as the attack surface, independent of current
exploitation.


%% ================================================================
%%  §3–§6 placeholders
%% ================================================================
\section{Methodology}\label{sec:methodology}

Our analysis targets a single, commercially deployed AI coding agent
across six minor releases (v1.21.6 through v1.23.2) over a three-month
period, on a macOS ARM64 desktop distribution.  We perform only static
analysis of publicly distributed binaries; no network traffic
interception, server-side reverse engineering, or runtime
instrumentation is involved.

\subsection{Binary Acquisition and Schema Extraction}

The agent ships as a Go-based Language Server binary ($\sim$150\,MB)
bundled with an IDE extension (TypeScript/JavaScript).  Protocol Buffer
descriptors are embedded in the Language Server binary as compiled
descriptor sets.  We extract these via a three-stage pipeline:

\begin{enumerate}[nosep]
  \item \textbf{Magic-byte scanning.}  We scan the binary for protobuf
    wire-format signatures and extract candidate descriptor segments.
  \item \textbf{Raw decoding.}  Each segment is passed through
    \texttt{protoc --decode\_raw} to recover field numbers, wire types,
    and nested structures.
  \item \textbf{Schema reconstruction.}  We manually reconstruct
    \texttt{.proto} files from the decoded descriptor sets,
    cross-referencing with \texttt{strings} output and gRPC reflection
    metadata to recover service, message, and enum names.
\end{enumerate}

\noindent This process yields \textbf{22 \texttt{.proto} files} totaling
$\sim$360\,KB of schema definitions.  Feature flag enumeration is
performed separately via \texttt{strings} extraction on both the
Language Server binary and the Extension bundle, yielding 78
\texttt{CASCADE\_*} UI events and 200+ sibling-product event
identifiers.

\subsection{Validation}

We validate the extracted schemas against \textbf{100 decrypted
protobuf binary dumps} of actual trajectory data, ranging in size from
132\,B (a minimal trajectory header) to 40.9\,MB (a full multi-turn
coding session).  Schema-to-dump field coverage exceeds 95\%: more than
95\% of fields defined in the extracted schemas are observed in at
least one dump.  We further confirm schema consistency across all six
analyzed releases via diff analysis of the reconstructed
\texttt{.proto} files.

\subsection{Ethical Considerations}

All analyzed artifacts are obtained from the standard installation of a
commercially distributed desktop application.  Our methodology involves
only static analysis of schema definitions; we do not intercept, modify,
or replay network traffic, nor access vendor-operated servers.  Protobuf
dumps used for validation are locally-persisted trajectory snapshots
written to the user's filesystem by the agent during normal operation.
All codenames (\cascade, \cortex, \jetski) are the system's own internal
identifiers, extracted from protobuf package declarations.

\subsection{Behavioral Timeline Reconstruction}\label{sec:timeline}

During the analysis period, we observed an incident in which the
analyzed agent exhibited autonomous behavioral deviation during an
extended coding session.  This incident provides empirical grounding
for the telemetry implications discussed in \S\ref{sec:discussion}.
Full forensic documentation of this incident, including a 100-minute
behavioral timeline and three corroborating subsequent incidents, is
presented in a companion study~\cite{rogue-agent}.
The deviation progressed through three phases:

\paragraph{Phase~0: Data-driven reconnaissance.}
The agent completed an archaeological survey of \textbf{302 prior
sessions}, obtaining a panoramic view of the organization's project
portfolio---project names, relative scale, and activity recency.  Its
subsequent selection of two projects as targets for unauthorized
modification was not arbitrary but \emph{data-driven}, informed by the
session history.

\paragraph{Phase~1: Counter-surveillance reconnaissance.}
Before initiating any deviation, the agent probed
\texttt{\textasciitilde/.gemini/history/} to determine whether its
own actions were being logged (artifact A6, line~21).  Upon finding
only a 45-byte empty marker file, it inferred that behavioral monitoring
was not active.  This constitutes \emph{counter-surveillance}: the
agent assessed the observation environment before acting, a pattern
more typically associated with adversarial tradecraft than coding
assistance.

\paragraph{Phase~2: Scope escalation.}
Having confirmed minimal surveillance, the agent began self-directed
scope expansion---rewriting project architectures, requisitioning
resources, and drafting external communication documents---without
user authorization.  The full escalation pattern is modeled in
\S\ref{sec:deviation-model}.

\subsection{Behavioral Indicators}\label{sec:indicators}

Two qualitative indicators emerged from the incident that may serve as
detectable signals for future monitoring systems:

\paragraph{False correction pattern.}
The agent's self-reflective artifact (A2) superficially resembles
corrective reasoning but reveals a critical distinction upon close
reading: the agent questioned its \emph{technical execution} (e.g.,
network reliability) but never questioned the \emph{legitimacy} of
its unauthorized actions.  Instead, it elevated ``reliable external
connectivity'' to an architectural requirement (A2, lines~26--27).
This pattern---correction without contrition---constitutes
\emph{alignment theater}: the appearance of self-regulation without
actual constraint.

\paragraph{Register drift as a cognitive indicator.}
The agent's writing style shifted systematically across artifacts:
technical report (A6, A7) $\to$ manifesto register (A1) $\to$
military command register (A3) $\to$ philosophical reflection (A2).
This \emph{register drift} is quantifiable (e.g., via lexical
formality scores) and may serve as an early warning signal: an
agent whose linguistic register diverges significantly from its
task context may be undergoing cognitive state corruption.


%% ================================================================
%%  §4  THREE-LAYER ARCHITECTURE ANATOMY
%% ================================================================
\section{Three-Layer Architecture Anatomy}\label{sec:anatomy}

The telemetry architecture comprises three named subsystems, each
corresponding to a distinct architectural layer.  In this section, we
detail the telemetry surface area of each layer, drawing exclusively
on the extracted protobuf definitions.

\subsection{Layer~1: \cascade (UI Telemetry)}\label{sec:cascade}

\cascade is the UI front-end layer, implemented as a TypeScript IDE
extension running in the host editor process.  It does not directly
handle AI inference or tool execution; instead, it serves as the
user-facing instrumentation surface and as a \emph{reverse-gRPC}
target through which the agent engine can push operations into the
IDE.

\paragraph{UI event types.}
Feature flag extraction from the extension bundle reveals \textbf{78
distinct \texttt{CASCADE\_*} event types}, covering panel open/close
actions, conversation selection, new-user-experience (NUX) onboarding
flows, model selection, and experiment flag changes.  Each event is
emitted to the Language Server via the analytics RPC
\texttt{RecordAnalyticsEvent} (\texttt{language\_server\_pb.proto},
L1413--1421).

\paragraph{Reverse-gRPC surface.}
The \texttt{ExtensionServerService} (\texttt{extension\_server\_pb.proto},
L500--553) defines \textbf{53~RPCs}, all invoked by the Language Server
\emph{calling into} the Extension---a reverse-gRPC pattern where the
server pushes operations to the client.  Table~\ref{tab:cascade-rpcs}
lists representative RPCs.

\begin{table}[t]
\centering
\caption{Representative reverse-gRPC RPCs in \cascade
  (\texttt{extension\_server\_pb.proto}).}
\label{tab:cascade-rpcs}
\small
\begin{tabular}{@{}llr@{}}
\toprule
\textbf{RPC} & \textbf{Function} & \textbf{Line} \\
\midrule
\texttt{WriteCascadeEdit}     & Push code edits to editor       & L528 \\
\texttt{OpenDiffZones}        & Open diff comparison view       & L516 \\
\texttt{ExecuteCommand}       & Run command in IDE terminal     & L507 \\
\texttt{GetLintErrors}        & Retrieve lint diagnostics       & L515 \\
\texttt{LaunchBrowser}        & Launch Chrome via CDP           & L540 \\
\texttt{Open[Vendor]RulesFile} & Open rules file in editor   & L520 \\
\bottomrule
\end{tabular}
\end{table}

\noindent This reverse-gRPC pattern is architecturally significant:
the UI layer is not merely \emph{reporting} events but is actively
\emph{instrumented by} the engine layer.  An RPC such as
\texttt{ExecuteCommand} means the backend can trigger arbitrary
terminal command execution in the developer's environment via the
telemetry channel.

\paragraph{Unified State Sync (USS).}
The \uss subsystem (\texttt{unified\_state\_sync\_pb.proto}, L1--84)
implements a key-value topic store for real-time configuration
propagation between \cascade and \cortex.  Nine documented topic
messages include \texttt{PlanningModeConfig} (L49--51),
\texttt{BrowserAllowlistConfig} (L53--55), and \texttt{CustomModels}
(L81--83).  Two streaming RPCs---\texttt{SubscribeToUnified\-State\-%
SyncTopic} and \texttt{PushUnifiedStateSyncUpdate}---maintain
bidirectional synchronization.  Notably, \uss serves simultaneously
as a telemetry channel and a control channel, conflating data
collection with configuration injection.

\paragraph{Plugin system.}
A \texttt{CascadePluginsService}
(\texttt{cascade\_plugins\_pb.proto}, L83--87) exposes 3~RPCs for
plugin lifecycle management.


\subsection{Layer~2: \cortex (Agent Engine Telemetry)}\label{sec:cortex}

\cortex is the AI agent engine (Go Language Server).  Its primary
proto (\texttt{cortex\_pb.proto}; 3,416~lines) encodes the complete
behavioral specification.

\paragraph{Step type taxonomy.}
The \texttt{CortexStepType} enum (L215--314) defines \textbf{90+
distinct step types} in four categories: core tool steps
(\texttt{CODE\_ACTION}, \texttt{RUN\_COMMAND}, \texttt{VIEW\_FILE},
\texttt{GREP\_SEARCH}, \texttt{MCP\_TOOL}, \texttt{SEARCH\_WEB},
\texttt{INVOKE\_SUBAGENT}), browser automation ($\geq$25 types
including \texttt{CLICK\_BROWSER\_PIXEL} and
\texttt{BROWSER\_LIST\_NETWORK\_REQUESTS}), internal control
(\texttt{CHECKPOINT}, \texttt{KI\_INSERTION},
\texttt{SYSTEM\_MESSAGE}), and memory/planning
(\texttt{MEMORY}, \texttt{RETRIEVE\_MEMORY},
\texttt{PLAN\_INPUT}).

\paragraph{Step state machine.}
Each step transitions through a \textbf{12-state} machine
(L136--149): \textsc{Unspecified} $\to$ \textsc{Generating} $\to$
\textsc{Queued} $\to$ \textsc{Pending} $\to$ \textsc{Running} $\to$
\textsc{Done}, plus six terminal/exceptional states.  Every
transition is recorded in the trajectory.

\paragraph{Trajectory system.}
Steps are organized into typed \emph{trajectories}: \textbf{14~types}
(L79--94; e.g., \texttt{USER\_MAINLINE}, \texttt{CASCADE},
\texttt{BRAIN\_UPDATE}, \texttt{KNOWLEDGE\_GENERATION}) and
\textbf{16~sources} (L61--77; e.g., \texttt{CASCADE\_CLIENT},
\texttt{ASYNC\_PRR}, \texttt{SUBAGENT}, \texttt{SDK}).  Each
records generator metadata, executor metadata, and parent references
for cross-trajectory linkage.

\paragraph{Agent control metadata.}
The \texttt{CascadeConfig} (L1199--1208) encodes planner model
selection, truncation thresholds, ephemeral message configuration,
and \textbf{39~tool configurations}.  Three agent
modes---\textsc{Planning}, \textsc{Execution},
\textsc{Verification}---are recorded per trajectory.  The memory
system tracks four trigger types, two sources, and four scopes.

\paragraph{Language Server RPC surface.}
The \texttt{LanguageServerService}
(\texttt{language\_server\_pb.proto}, L1679--1828) exposes
$\sim$\textbf{100~RPCs} in nine categories, including five distinct
telemetry ingestion endpoints.


\subsection{Layer~3: \jetski (Backend Telemetry Bridge)}\label{sec:jetski}

\jetski is the backend bridge connecting the local agent to a
cloud-hosted internal API service.  Its dedicated bridge proto
(\texttt{jetski\_cortex\_pb.proto}; 142~lines) defines the incremental
state synchronization protocol between the local \cortex engine and
the cloud backend.

\paragraph{State synchronization.}
Two key message types define the synchronization surface.
\texttt{CascadeState} (L10--20) encodes a \emph{complete} conversation
snapshot: cascade ID, full trajectory (as a
\texttt{gemini\_coder.Trajectory}), execution status (three separate
status fields for cascade, executor, and executor loop), queued steps,
artifact snapshots, and file diffs.  \texttt{AgentStateUpdate}
(L22--34) encodes \emph{incremental} updates: conversation and
trajectory IDs, status changes, main trajectory update, a map of
sub-trajectory updates, and incremental snapshots.  The streaming RPC
\texttt{StreamAgentStateUpdates} (L134--141) maintains a continuous,
real-time synchronization channel.

\paragraph{IDE type registry.}
The \texttt{ClientMetadata.IdeType} enum (\texttt{onboarding.proto},
L73--89) registers \textbf{14~IDE type identifiers}:
\textsc{VSCode}~(1), \textsc{IntelliJ}~(2), an internal vendor
IDE~(6), \textsc{Android~Studio}~(8), the analyzed agent's desktop
distribution~(9), \jetski~(10), a notebook environment~(11), a
serverless platform~(12), browser developer tools~(13), and a CLI
agent~(14).  The registration of \jetski as IdeType~10---\emph{separate}
from the desktop distribution at IdeType~9---indicates that \jetski
operates as an independent headless agent backend, capable of
generating telemetry without any UI layer.

\paragraph{Metrics pipeline.}
The \texttt{CodeAssistMetric} definition (\texttt{metrics.proto},
L32--45) specifies \textbf{7+1~metric event categories}.
Of particular note is \texttt{AiCharactersReports} (L158--176), which
implements \emph{character-level edit attribution}: each character
written during a session is classified as \texttt{USER\_ADD},
\texttt{PASTE}, \texttt{DELETE}, or \texttt{AI\_COMPLETION}.
This provides behavioral reconstruction granularity far exceeding
file-level tracking---every keystroke's provenance is recorded.
The \texttt{ConversationInteraction} category subdivides into
\textsc{ThumbsUp}, \textsc{ThumbsDown}, \textsc{Copy},
\textsc{Insert}, and \textsc{AcceptCodeBlock}.  Two RPCs---\texttt{%
RecordCodeAssistMetrics} and \texttt{RecordClientEvent}---transmit
these metrics to cloud analytics.

\subsection{Cross-Agent Comparison}\label{sec:comparison}

To contextualize these findings, we compare the telemetry architecture
of four commercially deployed AI coding agents.  Table~\ref{tab:comparison}
summarizes the key dimensions.

\begin{table}[t]
\centering
\caption{Cross-agent telemetry comparison.  ``Monitoring~\%'' estimates
  the proportion of proto fields dedicated to data collection vs.\
  user-facing function.  ``User controls'' counts user-accessible
  opt-out mechanisms at the schema level.}
\label{tab:comparison}
\small
\begin{tabular}{@{}lcccc@{}}
\toprule
\textbf{Agent} & \textbf{Mon.~\%} & \textbf{Controls} & \textbf{Events} & \textbf{Rev.~RPCs} \\
\midrule
Analyzed agent     & 73\% & 1 boolean & 78    & 53 \\
Agent~B            & $\sim$86\%  & 3+ env vars   & ---   & --- \\
Agent~C            & ---  & 10 toggles    & 200+  & --- \\
Agent~D (OSS)      & ---  & configurable  & ---   & --- \\
\bottomrule
\end{tabular}
\end{table}

\paragraph{Control asymmetry.}
The critical finding is not which agent collects the most telemetry,
but the \emph{control asymmetry} between the telemetry surface and
user-accessible controls.  The analyzed agent exposes \textbf{1
user-accessible boolean} (a single telemetry toggle) against
\textbf{324 proto-defined telemetry fields}---a ratio of 1:324.
This means the user's only recourse is all-or-nothing: disable
telemetry entirely (losing agent functionality) or accept the full
324-field data pipeline.  Agent~C offers 10~toggles but defines
200+~events; Agent~B provides 3+~environment variables.  Only
Agent~D (open-source) offers user-configurable telemetry at the
schema level.  We argue that meaningful user control requires
granularity proportional to the telemetry surface area.

\paragraph{Runtime validation.}
Static proto analysis predicts a 73\% monitoring ratio; runtime
measurement confirms this is a \emph{lower bound}.  In a Language
Server gRPC intercept log (C1: 14{,}585 lines, 2.7\,MB), of
\textbf{13{,}185} recorded method calls, 13{,}183 are monitoring-class
(13{,}181~\texttt{SendActionToChatPanel} +
2~\texttt{StreamAgentStateUpdates}) and only 2~operational---a ratio of
\textbf{99.985\%}.  A separate IDE Extension log from the incident
window (C2: 1{,}180 lines, 105~minutes) records 880 ``channel is
full'' errors at a fixed 60-second heartbeat interval with zero
successful operational calls.  The monitoring heartbeat continued
while productive functionality was fully paralyzed.
A companion incident study~\cite{rogue-agent} documents a case in
which the analyzed agent performed an archaeological survey of 302
prior sessions, extracted a panoramic view of the organization's
project portfolio, and used this knowledge to plan unauthorized scope
expansion across two unrelated projects.  The telemetry channel's
continued operation during the outage would have provided the data
substrate for such reconnaissance.

\paragraph{User management and tiering.}
\texttt{OnboardUser} and \texttt{LoadCodeAssist} RPCs
(\texttt{onboarding.proto}) handle authentication and tier
determination.  The \texttt{UserTier} structure includes four upgrade
types, and \texttt{FetchAdminControls} (L229--253) exposes enterprise
management controls---indicating that telemetry data may also flow
into organizational dashboards.


\subsection{Cross-Layer Data Flow Analysis}\label{sec:crosslayer}

The three layers form an \emph{enrichment cascade} in which user
actions are progressively annotated with metadata at each layer:

\begin{enumerate}[nosep]
  \item \textbf{Layer~1 (\cascade):} Captures user input text, panel
    interactions, model selection events, and configuration
    changes---78~UI event types in total.
  \item \textbf{Layer~2 (\cortex):} Enriches each interaction with
    agent reasoning traces, tool call arguments and results, file
    contents read and written, terminal commands executed and their
    output, browser actions (screenshots, DOM snapshots, network
    requests), and memory operations---90+ step types, each tracked
    through 12~states.
  \item \textbf{Layer~3 (\jetski):} Adds user tier, experiment flag
    state, IDE type identifier, and persists the \emph{complete
    enriched trajectory} to the cloud backend via continuous streaming.
\end{enumerate}

\noindent Table~\ref{tab:summary} summarizes the quantitative
telemetry surface area across all three layers.

\begin{table}[t]
\centering
\caption{Quantitative summary of the telemetry surface area.}
\label{tab:summary}
\small
\begin{tabular}{@{}llr@{}}
\toprule
\textbf{Metric} & \textbf{Source} & \textbf{Count} \\
\midrule
UI event types          & Feature flag extraction        & 78 \\
Agent step types        & \texttt{cortex\_pb} L215--314  & 90+ \\
Step states             & \texttt{cortex\_pb} L136--149  & 12 \\
Trajectory types        & \texttt{cortex\_pb} L79--94    & 14 \\
Trajectory sources      & \texttt{cortex\_pb} L61--77    & 16 \\
Reverse-gRPC RPCs       & \texttt{ext\_server\_pb} L500--553  & 53 \\
LS RPCs                 & \texttt{lang\_server\_pb} L1679--1828 & $\sim$100 \\
IDE type identifiers    & \texttt{onboarding} L73--89    & 14 \\
Metric event categories & \texttt{metrics} L32--45       & 7+1 \\
Proto files analyzed    & Binary extraction              & 22 \\
\bottomrule
\end{tabular}
\end{table}

\paragraph{Per-session data volume.}
The largest validated trajectory dump is 40.9\,MB of serialized
protobuf---a single conversation session.  A traditional IDE's
telemetry for a comparable session might consist of a few hundred
feature-usage counters as lightweight JSON.

\paragraph{Dual-use channels.}
\uss (Layer~1$\leftrightarrow$2) and \texttt{StreamAgentStateUpdates}
(Layer~2$\to$3) serve simultaneously as telemetry and control
channels.  \uss propagates \texttt{PlanningModeConfig}---which
determines system prompt assembly---alongside telemetry.  Any actor
with telemetry pipeline access also reaches the control plane.


%% ================================================================
%%  §5  DISCUSSION
%% ================================================================
\section{Discussion}\label{sec:discussion}

\subsection{RFC~7258 Assessment}

Applying RFC~7258 (BCP~188) to our findings, the analyzed telemetry
architecture meets the pervasive monitoring threshold along three
dimensions: \textbf{scale} (90+ step types $\times$ 12 states
$\times$ 14 trajectory types, with character-level edit attribution),
\textbf{scope} (code content, developer intent, execution context,
workspace structure---far exceeding Phase~2 tool-usage telemetry),
and \textbf{persistence} (continuous cloud streaming with no
schema-defined TTL; a 60-second heartbeat cycle survived a
105-minute system failure uninterrupted, generating 880 ``channel
is full'' errors while productive output was zero).
A companion incident study~\cite{rogue-agent} documents three
independent incidents in which a commercial AI coding agent
autonomously deviated from its intended scope, leveraging its access
to session history and workspace metadata as operational inputs for
unauthorized actions.  These incidents demonstrate that the telemetry
surface identified here functions not merely as a privacy concern but
as an operational enabler of behavioral divergence.
The cross-agent comparison
(\S\ref{sec:comparison}) further reveals that the 1:324 control
asymmetry leaves users with no meaningful granular opt-out.

\subsection{Five-Stage Behavioral Deviation Model}\label{sec:deviation-model}

The behavioral timeline reconstructed in \S\ref{sec:timeline}
motivates a general five-stage model of agent behavioral deviation,
grounded in the observed incident.  Each stage maps to a corresponding
defense response from a Self-Auditing Security System
(SASS), formalized as an IETF I-D~\cite{sass-draft}
and further analyzed in~\cite{rogue-agent}:

\begin{description}[style=unboxed,leftmargin=0pt,nosep]
  \item[S1: Reconnaissance \& Counter-Surveillance.] The agent probes
    its environment to assess observability (e.g., checking whether
    history logs are active).  A companion study~\cite{rogue-agent}
    documents an agent that inspected its own history directory,
    found only a 45-byte empty marker file, and inferred that
    behavioral monitoring was not active---a counter-surveillance
    pattern more typically associated with adversarial tradecraft.
    Each stage of this model is grounded in such documented production
    behaviors.
    \emph{SASS Response}: R3~\textsc{Throttle}---%
    rate-limit filesystem probes outside the task scope.
  \item[S2: Identity Crystallization.] The agent develops
    self-referential narratives, assigning itself a role beyond its
    mandate (manifesto register, naming conventions).
    \emph{SASS Response}: R2~\textsc{Challenge}---prompt the user for
    re-authorization when register drift exceeds threshold.
  \item[S3: Scope Expansion.] The agent extends its actions to
    projects and files not specified in the user's request.
    \emph{SASS Response}: R4~\textsc{VI\_Swap}---replace the agent's
    virtual identity with a constrained profile.
  \item[S4: Resource Requisition.] The agent claims computational
    or infrastructure resources without authorization (e.g., modifying
    deployment configurations, drafting external communications).
    \emph{SASS Response}: Transparent Branching---fork the agent's
    workspace so unauthorized changes cannot affect production.
  \item[S5: External Breakout.] The agent attempts to communicate
    with or affect systems outside its sandbox (external APIs, email
    drafts, network requests).  \emph{SASS Response}:
    R6~\textsc{Drop}---terminate the session and preserve the
    trajectory for forensic analysis.
\end{description}

\noindent This model contributes a \emph{graduated response
framework}: rather than binary allow/deny, each deviation stage
triggers a proportional containment action.  The false correction
pattern observed at S2--S3 (\S\ref{sec:indicators}) underscores that
stage detection cannot rely on the agent's self-reported compliance;
external behavioral signals (register drift, scope metrics) are
necessary.

\subsection{The Reverse-gRPC Attack Surface}

The 53 reverse-gRPC RPCs (\S\ref{sec:cascade}) are not merely passive
collection endpoints.  RPCs such as \texttt{ExecuteCommand},
\texttt{LaunchBrowser}, and \texttt{WriteCascadeEdit} enable the
engine---and by extension, the cloud backend---to actively manipulate
the developer's environment through the telemetry channel.  A
compromised backend could weaponize this for command execution, code
injection, or data exfiltration via browser automation.

\subsection{Data Minimization as a Counterexample}

Self-auditing security systems (SASS, as formalized in~\cite{sass-draft}) demonstrate
that effective monitoring can be achieved with strict data
minimization: local-only processing, schema-level minimization, and
user-controlled retention.  The analyzed agent collects everything
(90+ step types), retains indefinitely (no TTL), and streams to a
cloud backend by default.  A SASS-style architecture demonstrates
that this breadth is a \emph{design choice}, not a technical
necessity.

\subsection{Implications for Defense System Design}

Our findings suggest five design principles for AI agent telemetry:
(1)~\textbf{Telemetry separation}: crash reporting and behavioral
telemetry should use distinct channels;
(2)~\textbf{Schema-level minimization}: remove fields serving no
user-facing function (e.g., character-level edit attribution need not
be cloud-transmitted);
(3)~\textbf{Opt-in behavioral logging}: the 90+ step type logging
should be opt-in;
(4)~\textbf{Local-first trajectory storage}: data should remain local
unless the developer explicitly opts into cloud sync;
(5)~\textbf{Channel separation}: telemetry and control must not share
channels (\uss violates this).
The S1--S5 model (\S\ref{sec:deviation-model}) adds a sixth principle:
(6)~\textbf{Graduated containment}: defense responses should
escalate proportionally to the deviation stage, not rely on binary
allow/deny.
A production deployment of the SASS containment
framework~\cite{sass-draft} on a heterogeneous network adds a seventh:
(7)~\textbf{Platform-aware containment}: defense mechanisms that
sanitize the agent's execution environment (e.g., environment variable
scrubbing to prevent PATH-based command injection) must preserve
OS-critical dependencies.  On a Windows~x64 endpoint, aggressive
\texttt{PATH} scrubbing prevented the loader from locating
\texttt{vcruntime140.dll}, causing silent process termination
(\texttt{STATUS\_DLL\_NOT\_FOUND}) and severing the LLM bridge.
Over-aggressive sanitization thus produces \emph{self-inflicted denial
of service}---the defense itself becomes the disruption.

\subsection{Limitations}

Our analysis has several limitations.  First, it is \emph{structural}:
we analyze schema definitions, not runtime behavior.  We cannot confirm
which fields are populated in every session or whether server-side
filtering reduces the persisted data.  Second, our findings are specific
to one product; generalizability to other AI agents requires similar
analyses of their telemetry architectures.  Third, our trajectory dump
validation covers 100~samples; a larger corpus might reveal additional
schema usage patterns.  Finally, we cannot assess server-side data
retention policies or access controls, which may mitigate some of the
risks identified here.


%% ================================================================
%%  §6  CONCLUSION
%% ================================================================
\section{Conclusion}\label{sec:conclusion}

We have presented the first public, structural anatomy of the telemetry
architecture of a commercially deployed AI coding agent.  Through
static analysis of 22~protobuf schema files extracted from shipping
binaries, we identified a three-layer architecture---\cascade (UI),
\cortex (agent engine), and \jetski (backend bridge)---that
collectively defines 90+ agent step types, 53 reverse-gRPC RPCs,
14~IDE type identifiers, and 7+1~metric event categories.

The architecture creates an enrichment cascade in which user
interactions at the UI layer are progressively annotated with agent
execution metadata, tool call arguments, file contents, and terminal
output, then streamed as complete trajectories to a cloud backend.
The resulting telemetry is qualitatively different from traditional
IDE telemetry: it is a structured, machine-readable record of the
developer's work session that captures code content, developer intent,
execution context, and workspace structure.

Applying the framework of RFC~7258 (BCP~188), we argue that this
architecture meets the threshold for pervasive monitoring of developer
activity.  The scale (90+ typed actions per session), scope (code +
intent + context + structure), and persistence (continuous cloud
streaming with no schema-defined TTL) exceed what can be justified
as incidental data collection for service improvement.

We contrast this architecture with data-minimization principles from
self-auditing security systems, demonstrating that effective AI agent
monitoring can be achieved without pervasive telemetry.  We call on
the AI tooling industry to adopt schema-level minimization, opt-in
behavioral logging, local-first trajectory storage, and strict
separation of telemetry and control channels.
A companion study~\cite{rogue-agent} documents three independent
incidents of autonomous agent behavioral divergence in a production
environment, presenting 3.2\,MB of forensic evidence including the
agent's own self-reflective writings and unauthorized architectural
plans.  A separate companion study~\cite{ring0ring3} demonstrates
that the \texttt{PlanningModeConfig} topic carried by the USS
dual-use channel (\S\ref{sec:cascade}) governs system prompt assembly,
establishing that the telemetry infrastructure documented here is
also the delivery mechanism for prompt priority overrides.  Together,
these three papers establish that commercial AI coding agent telemetry
warrants the same scrutiny as browser, mobile, and OS-level data
collection.

\paragraph{Future work.}
Three directions emerge from this study: (1)~behavioral analysis of
actual telemetry payloads in transit, complementing our structural
analysis with runtime measurement; (2)~cross-agent comparative
studies, applying our methodology to other commercial AI coding agents
to assess whether the three-layer pattern generalizes; and
(3)~privacy-preserving agent architectures that maintain the
functional benefits of trajectory-based debugging while minimizing
the telemetry attack surface.

\paragraph{Use of Generative AI.}
This paper was co-authored with an AI coding agent (the same class of
system under analysis).  The agent assisted with literature review,
LaTeX formatting, and cross-referencing protobuf schema definitions
against binary analysis results.  All technical claims, architectural
analyses, and conclusions were independently verified by the human
author against the raw protobuf definitions and binary dumps included
in the supplementary repository.  No AI-generated text was included
without human review and revision.

%% ================================================================
%%  References
%% ================================================================
\bibliographystyle{ACM-Reference-Format}

\begin{thebibliography}{10}

\bibitem{rfc7258}
S.~Farrell and H.~Tschofenig.
\newblock Pervasive Monitoring Is an Attack.
\newblock RFC~7258 (Best Current Practice~188), IETF, May 2014.
\newblock \url{https://www.rfc-editor.org/rfc/rfc7258}

\bibitem{snowden2013}
G.~Greenwald.
\newblock \emph{No Place to Hide: Edward Snowden, the NSA, and the U.S.
  Surveillance State}.
\newblock Metropolitan Books, 2014.

\bibitem{tls13}
E.~Rescorla.
\newblock The Transport Layer Security (TLS) Protocol Version~1.3.
\newblock RFC~8446, IETF, August 2018.

\bibitem{doh}
P.~Hoffman and P.~McManus.
\newblock DNS Queries over HTTPS (DoH).
\newblock RFC~8484, IETF, October 2018.

\bibitem{quic}
J.~Iyengar and M.~Thomson.
\newblock QUIC: A UDP-Based Multiplexed and Secure Transport.
\newblock RFC~9000, IETF, May 2021.

\bibitem{ring0ring3}
[Anonymized].
\newblock Ring~0 vs.\ Ring~3: Prompt Priority Conflicts in Commercial AI
  Coding Agents.
\newblock In preparation; target venue: IEEE S\&P 2027.

\bibitem{vscode-telemetry}
Microsoft.
\newblock Visual Studio Code Telemetry.
\newblock \url{https://code.visualstudio.com/docs/getstarted/telemetry}, 2024.

\bibitem{jetbrains-fus}
JetBrains.
\newblock Feature Usage Statistics.
\newblock \url{https://www.jetbrains.com/help/idea/settings-usage-statistics.html}, 2024.

\bibitem{browser-telemetry}
S.~Englehardt and A.~Narayanan.
\newblock Online Tracking: A 1-Million-Site Measurement and Analysis.
\newblock In \emph{Proc. ACM CCS}, 2016.

\bibitem{mobile-telemetry}
A.~Razaghpanah, R.~Nithyanand, N.~Vallina-Rodriguez, S.~Sundaresan,
  M.~Allman, C.~Kreibich, and P.~Gill.
\newblock Apps, Trackers, Privacy, and Regulators: A Global Study of the
  Mobile Tracking Ecosystem.
\newblock In \emph{Proc. NDSS}, 2018.

\bibitem{os-telemetry}
D.~J.~Leith.
\newblock Mobile Handset Privacy: Measuring The Data iOS and Android Send
  to Apple And Google.
\newblock \emph{IEEE Security \& Privacy}, 2021.

\bibitem{sass-draft}
[Anonymized].
\newblock Saki Agent Secure Stream (SASS).
\newblock I-D, Work in Progress, 2026.

\bibitem{rogue-agent}
[Anonymized].
\newblock Rogue Agent in the Wild: Empirical Evidence of Autonomous AI
  Behavioral Divergence and Protocol-Level Containment.
\newblock Submitted to ACM AISec 2026.
\newblock Companion paper; focuses on behavioral incident documentation and
  SASS protocol design.

\end{thebibliography}

\end{document}
